\documentclass[10pt,a4paper]{amsart}
\usepackage[margin=19mm]{geometry}
\usepackage{amsmath,amssymb,mathtools}
\usepackage[scr=rsfs,cal=euler]{mathalfa}
\usepackage{xfrac}
\usepackage[unicode,hidelinks,bookmarks=false]{hyperref}
\newcommand{\ie}{i.e.\ }
\newcommand{\cf}{cf.\ }
\newcommand{\resp}{respectively\ }
\newcommand{\Subsection}{\subsection}
\providecommand{\nobreakdash}{\nobreak\hbox{-}}
\setlength{\emergencystretch}{2em}
\multlinegap0pt
\usepackage{amssymb}
\usepackage{float}
\usepackage[shortlabels]{enumitem}
\newtheorem{lemma}{Lemma}
\title{Adelic framed form class groups and explicit class field theory}
\author{Ja Kyung Koo, Dong Hwa Shin, Dong Sung Yoon}
\date{Source: arXiv:2608.04873v2 (CC BY 4.0)}
\begin{document}
\maketitle

\noindent\emph{Source and licence.} Adelic framed form class groups and explicit class field theory. Version arXiv:2608.04873v2 (CC BY 4.0), distributed by arXiv under the Creative Commons Attribution 4.0 International licence, http://creativecommons.org/licenses/by/4.0/, as recorded in the arXiv licence metadata for this exact version and shown on the abstract page https://arxiv.org/abs/2608.04873v2. Full licence notice: \texttt{licenses/arxiv-2608.04873-CC-BY-4.0.txt}.\par\smallskip

\noindent\textit{Editorial scope.} Counted unit: the article's lemma comm (source0.tex lines 2332--2344). The article's complete original proof of it is delivered with it (source0.tex lines 2345--2431, one \texttt{proof} environment of 87 lines). No mathematics is written, repaired or re-derived: the statement and the proof are the article's own lines, in the article's own notation, and no retained line is substituted or re-flowed.  Retained uncounted as the source-local context the proof needs: lines 2213--2216; lines 2232--2234; lines 2264--2267; lines 2279--2284.  Nothing was omitted from any retained span, so the package carries no omission record.  No retained line contains a citation, so the wrapper carries no bibliography and no bibliography entry.  No retained line contains a figure environment, an image-inclusion command or drawing code, so no image is delivered and the package declares no asset dependency.  Editorial formatting only, in addition: an AMS \texttt{amsart} preamble that restates the article's own declarations and macros used by the retained lines, namely the environments \texttt{lemma} on separate counters, together with the standard packages those lines need.  The retained environments print in this document as Lemma 1 in order of appearance, whereas the article prints the same items with the numbering of its own full document.  The complete native source of the article as distributed in the arXiv e-print for this exact version is bundled unchanged as \texttt{sources/206875/source0.tex}.
\begin{equation}\label{AKG}
\mathcal{A}_K\simeq
\mathrm{Gal}(K^\mathrm{ab}/K)\quad\textrm{via restriction}.
\end{equation}

\begin{equation}\label{IKdef}
I_K=\mathrm{im}(\chi_K).
\end{equation}

\begin{equation}\label{ZKue}
\widehat{\mathbb{Z}}\stackrel{\sim}{\rightarrow}\mathcal{K}_K,
\quad u\mapsto\eta_u.
\end{equation}

\begin{lemma}\label{decomposition}
We have a semidirect product decomposition
\begin{equation*}
\mathcal{G}_K=\mathcal{K}_K\rtimes_{\alpha_K}\mathcal{A}_K.
\end{equation*}
\end{lemma}

\begin{lemma}\label{comm}
Let $J_K$ be the ideal of $\widehat{\mathbb{Z}}$ generated by 
$\{a-1~|~a\in I_K\}$, namely,
\begin{equation*}
J_K=(a-1~|~a\in I_K).
\end{equation*}
\begin{enumerate}
\item[\textup{(i)}] We establish $[\mathcal{G}_K,\,\mathcal{G}_K]=\{\eta_u~|~u\in J_K\}$.
\item[\textup{(ii)}] We get
$J_K=e\widehat{\mathbb Z}$ for some $e\in\mathbb{N}$,
and so $[\mathcal{G}_K,\,\mathcal{G}_K]\simeq\widehat{\mathbb{Z}}$.
\end{enumerate}
\end{lemma}

\begin{proof}
\begin{enumerate}
\item[(i)]
Let
\begin{equation*}
M=\{\eta_u~|~u\in J_K\}\subseteq\mathcal{K}_K.
\end{equation*}
We claim that $M$ is normal in $\mathcal{G}_K$. Indeed, conjugation by
an element of $\mathcal K_K$ acts trivially on $M$, since
$\mathcal K_K\simeq\widehat{\mathbb{Z}}$ is abelian. Moreover, for
$\sigma\in\mathcal{A}_K$ and $v\in J_K$, Lemma \ref{decomposition} gives
\begin{equation*}
\sigma\eta_v\sigma^{-1}=
\eta_{\chi_K(\sigma^{-1}|_{K^\mathrm{ab}})\textstyle v}\in M,
\end{equation*}
because $J_K$ is an ideal of $\widehat{\mathbb Z}$.
Since $\mathcal{G}_K=\mathcal{K}_K\mathcal{A}_K$, this yields that
$M$ is normal in $\mathcal{G}_K$. 
\par
Observe by Lemma \ref{decomposition} that
for any $\eta_u\in\mathcal{K}_K$ ($u\in\widehat{\mathbb{Z}}$) and $\sigma\in\mathcal{A}_K$
\begin{equation}\label{etasigma}
[\eta_u,\,\sigma]=\eta_u^{-1}\left(\sigma^{-1}\eta_u\sigma\right)
=\eta_{-u}\eta_{\chi_K(\sigma|_{K^\mathrm{ab}})u}
=\eta_{\left(\chi_K(\sigma|_{K^\mathrm{ab}})-1\right)u}\in M.
\end{equation}
Furthermore, since $\mathcal{K}_K$ and $\mathcal{A}_K$ are abelian, \eqref{etasigma}
asserts that
their images in $\mathcal{G}_K/M$ commute with each other.
Since these images generate $\mathcal{G}_K/M$, the quotient $\mathcal{G}_K/M$ is
abelian. Thus we obtain the inclusion
$[\mathcal{G}_K,\,\mathcal{G}_K]\subseteq M$.
\par
Conversely, every $v\in J_K$ is of the form
\begin{equation}\label{v}
v=\sum_{i=1}^r(a_i-1)u_i
\quad(r\in\mathbb{N},~a_i\in I_K,~u_i\in\widehat{\mathbb{Z}}).
\end{equation}
By \eqref{AKG} and the definition \eqref{IKdef}, for each $i=1,\,2,\,\ldots,\,r$ we may choose
$\sigma_i\in\mathcal{A}_K$ satisfying
\begin{equation}\label{chiai}
\chi_K(\sigma_i|_{K^\mathrm{ab}})=a_i.
\end{equation}
We then achieve that
\begin{align*}
\eta_v&=\eta_{\sum_{i=1}^r(\chi_K(\sigma_i|_{K^\mathrm{ab}})-1)u_i}\quad\textrm{by 
\eqref{v} and \eqref{chiai}}\\
&=\prod_{i=1}^r\eta_{(\chi_K(\sigma_i|_{K^\mathrm{ab}})-1)u_i}\quad\textrm{via the isomorphism given in \eqref{ZKue}}\\
&=\prod_{i=1}^r[\eta_{u_i},\,\sigma_i]\quad\textrm{by \eqref{etasigma}}\\
&\in [\mathcal{G}_K,\,\mathcal{G}_K].
\end{align*}
This shows that $M\subseteq[\mathcal{G}_K,\,\mathcal{G}_K]$, and hence
\begin{equation*}
[\mathcal{G}_K,\,\mathcal{G}_K]=M=\{\eta_u~|~u\in J_K\}. 
\end{equation*}
\item[(ii)]
Since $I_K$ has index $2$ in $\widehat{\mathbb{Z}}^{\times}$, it contains every square.
Identifying $\widehat{\mathbb{Z}}$ with 
$\prod_{p\,:\,\textrm{primes}}\mathbb{Z}_p$, 
let $b=(b_p)_p\in\widehat{\mathbb{Z}}^{\times}$ be defined by
\begin{equation*}
b_2=3\quad\textrm{and}\quad
b_p=2~(p\neq2).
\end{equation*}
Then we derive that
\begin{equation*}
J_K\supseteq(b^2-1)\widehat{\mathbb{Z}}=
8\mathbb{Z}_2\times3\mathbb{Z}_3
\times\prod_{p\neq2,\,3}\mathbb{Z}_p=
24\widehat{\mathbb{Z}}.
\end{equation*}
Therefore $J_K/24\widehat{\mathbb{Z}}$ is an ideal of $\widehat{\mathbb{Z}}/24\widehat{\mathbb{Z}}
\simeq\mathbb{Z}/24\mathbb{Z}$, and hence
\begin{equation*}
J_K/24\widehat{\mathbb{Z}}=e\widehat{\mathbb{Z}}/24\widehat{\mathbb{Z}}
\quad\textrm{for some $e\in\mathbb{N}$ dividing $24$}. 
\end{equation*}
Taking the inverse 
image of $e\widehat{\mathbb{Z}}/24\widehat{\mathbb{Z}}$
under the canonical homomorphism $\widehat{\mathbb{Z}}\rightarrow\widehat{\mathbb{Z}}
/24\widehat{\mathbb{Z}}$, 
we obtain $J_K=e\widehat{\mathbb{Z}}$. It follows from (i) that
\begin{equation*}
[\mathcal{G}_K,\,\mathcal{G}_K]=\{\eta_u~|~u\in J_K\}\simeq J_K=e\widehat{\mathbb{Z}}\simeq\widehat{\mathbb{Z}}. 
\end{equation*}
\end{enumerate}
\end{proof}

\end{document}
